% p208_2; source PDF page 208; manually verified editable table.
\begin{tikzpicture}[x=1cm,y=-1cm,inner sep=0pt,outer sep=0pt]
\path[use as bounding box] (0,0) rectangle (16.000,4.600);
\draw[line width=.4pt] (0.000,0.000) rectangle (4.000,0.850);
\node[align=center,text width=3.740cm,font={\linespread{1}\fontsize{10}{13.50}\selectfont\bfseries}] at (2.000,0.425) {事项};
\draw[line width=.4pt] (4.000,0.000) rectangle (16.000,0.850);
\node[align=center,text width=11.740cm,font={\linespread{1}\fontsize{10}{13.50}\selectfont\bfseries}] at (10.000,0.425) {程序};
\draw[line width=.4pt] (0.000,0.850) rectangle (4.000,2.100);
\node[align=left,text width=3.740cm,font={\linespread{1}\fontsize{10}{13.50}\selectfont}] at (2.000,1.475) {一般的减刑、假释};
\draw[line width=.4pt] (4.000,0.850) rectangle (16.000,2.100);
\node[align=left,text width=11.740cm,font={\linespread{1}\fontsize{10}{13.50}\selectfont}] at (10.000,1.475) {由执行机关向中级人民法院提出减刑、假释建议书，由中级人民法院作出裁定};
\draw[line width=.4pt] (0.000,2.100) rectangle (4.000,3.350);
\node[align=left,text width=3.740cm,font={\linespread{1}\fontsize{10}{13.50}\selectfont}] at (2.000,2.725) {无期徒刑的减刑、假释};
\draw[line width=.4pt] (4.000,2.100) rectangle (16.000,3.350);
\node[align=left,text width=11.740cm,font={\linespread{1}\fontsize{10}{13.50}\selectfont}] at (10.000,2.725) {由执行机关向高级人民法院提出减刑、假释建议书，由高级人民法院作出裁定};
\draw[line width=.4pt] (0.000,3.350) rectangle (4.000,4.600);
\node[align=left,text width=3.740cm,font={\linespread{1}\fontsize{10}{13.50}\selectfont}] at (2.000,3.975) {死缓的减刑};
\draw[line width=.4pt] (4.000,3.350) rectangle (16.000,4.600);
\node[align=left,text width=11.740cm,font={\linespread{1}\fontsize{10}{13.50}\selectfont}] at (10.000,3.975) {由执行机关向高级人民法院提出减刑建议书，由高级人民法院作出裁定};
\end{tikzpicture}
